\clearpage\SourceLeaf{517}
\renewcommand{\BookRunningTitle}{Auxiliary Table for Solving the Indeterminate Equation}
\pdfbookmark[0]{Auxiliary Table for Solving the Indeterminate Equation}{hulfstafel-ausschliessung}
\begin{center}
\textbf{AUXILIARY TABLE}\\[.7em]
FOR SOLVING THE INDETERMINATE EQUATION
\[
A=fxx+gyy
\]
BY THE METHOD OF EXCLUSION.
\end{center}
\BookRule{.18\textwidth}
It is assumed that a prime $p$ has been chosen as the exclusion number, and that none of the numbers $A$, $f$, $g$ is divisible by it. The table is also restricted to the two cases in which the value of the expression $\frac{A}{f}(\mathrm{mod}.p)$ is either a fixed quadratic residue (always $1$), or a fixed quadratic nonresidue modulo $p$. Finally, the table has been provided only for the case in which $fg$ is a quadratic residue of $p$, and the opposite case has been omitted entirely. Section Six of the \textit{Disquisitiones Arithmeticae} gives sufficient instruction on how to derive easily from the table what it does not contain directly.
\par\medskip\textit{Examples.} Let the equation to be solved be $21680143=xx+78yy$
\par\medskip\noindent 1) Exclusion number $=5$\hfill $fg=\mathrm{N},\quad \frac{A}{f}\equiv3\equiv2\pdot2^2$
\par\noindent From the table $1,4$ for $fg=R$ and hence $0,2,3$ for $fg=\mathrm{N}$,
and for the above case $0,1,4$ or excluded $5n\pm2$
\par\medskip\noindent 2) Exclusion number $=7$\hfill $fg=R,\quad \frac{A}{f}\equiv2\equiv1,3^2$
\par\noindent From the table $0,1,2,5,6$ For the present case $0,3,6,1,4$ and excluded $7n\pm2$
\par\medskip\noindent 3) Exclusion number $=11$\hfill $fg=R,\quad \frac{A}{f}\equiv1$
\par\noindent Thus we obtain $0,1,3,5,6,8,10$; excluded $11n\pm2,4$
\par\medskip\noindent 4) Exclusion number $=17$\hfill $fg=\mathrm{N},\quad \frac{A}{f}\equiv9\equiv1\pdot3^1$
\[
\begin{gathered}
0\pdot1\pdot3\pdot4\pdot6\\
0\pdot3\pdot8\pdot5\pdot1\\
0,2,3,4,6,7\\
\text{Exclusion number}\quad\pm1,5,8
\end{gathered}
\]
